\begin{figure}[htbp]
\centering
\begin{tikzpicture}[x=2.65cm,y=2.65cm,font=\small,
  primal/.style={lagoon,line width=1.05pt},
  dual/.style={ember,line width=.85pt,dash pattern=on 4pt off 2pt},
  axis/.style={-{Stealth[length=3pt]},line width=.45pt,black!65},
  marker/.style={fill=lagoon,draw=white,line width=.55pt}]
  \newcommand{\priceaxes}{%
    \fill[black!3] (0,0) rectangle (1,1);
    \draw[black!27,line width=.4pt] (0,1)--(1,1)--(1,0);
    \draw[axis] (-.19,0)--(1.22,0);
    \draw[axis] (0,-.2)--(0,1.24);
    \node[anchor=east,inner sep=2pt] at (-.015,1){\(1\)};
    \node[anchor=north,inner sep=2pt] at (1,-.105){\(1\)};
  }
  \begin{scope}
    \priceaxes
    \node at (.55,1.55){Original prices \(p=(0,0)\)};
    \draw[dual] (-.17,.010625)--(1.19,-.074375);
    \draw[primal] (-.01125,-.18)--(.075625,1.21);
    \path[marker] (.0625,1) circle[radius=2.3pt];
    \node[anchor=west,inner sep=2pt] at (.12,.87){\((1/16,1)\)};
    \node at (.53,-.40){\(z_1\) free,\quad \(z_2=+\)};
  \end{scope}
  \begin{scope}[xshift=7.2cm]
    \priceaxes
    \node at (.55,1.55){New prices \(p=(-4,0)\)};
    \draw[dual] (-.17,.17)--(.19,-.19);
    \draw[primal] (-.16,-.16)--(1.17,1.17);
    \path[marker] (1,1) circle[radius=2.3pt];
    \node[anchor=east,inner sep=2pt] at (.93,1.09){\((1,1)\)};
    \node at (.53,-.40){\(z_1=+,\quad z_2=+\)};
  \end{scope}
  \draw[-{Stealth[length=4pt]},black!65] (1.37,.55)--(2.15,.55);
\end{tikzpicture}
\caption{Making a small witness coordinate visible. The solid line is
the scaled primal space for \(C=\operatorname{span}\{(1/16,1)\}\), and
the dashed line is its orthogonal complement. The dot is the
nonnegative primal vector of sup norm one. Multiplying its first
coordinate by \(16\) and the corresponding dual coordinate by
\(1/16\) preserves orthogonality. At the new prices both coordinates
reach the visibility threshold \(1/2\), so both labels are forced.
Coordinates \(1\) and \(2\) are horizontal and vertical, respectively.}
\label{fig:scale}
\end{figure}
